
\documentclass{article}
\usepackage{colortbl}
\usepackage{makecell}
\usepackage{multirow}
\usepackage{supertabular}

\begin{document}

\newcounter{utterance}

\centering \large Interaction Transcript for game `cladder', experiment `full\_v1.5\_default', episode 8235 with qwen3.5{-}2b.
\vspace{24pt}

{ \footnotesize  \setcounter{utterance}{1}
\setlength{\tabcolsep}{0pt}
\begin{supertabular}{c@{$\;$}|p{.15\linewidth}@{}p{.15\linewidth}p{.15\linewidth}p{.15\linewidth}p{.15\linewidth}p{.15\linewidth}}
    \# & \multicolumn{2}{c}{Player} && \multicolumn{2}{c}{Game Master} \\
    \hline

    \theutterance \stepcounter{utterance}  
    & & & \multicolumn{4}{p{0.6\linewidth}}{
        \cellcolor[rgb]{0.9,0.9,0.9}{
            \makecell[{{p{\linewidth}}}]{
                \texttt{\tiny{[P1$\langle$GM]}}
                \texttt{You are an expert in causal inference. The following question is not a typical commonsense query, but rather a meticulously designed question created by a professor specializing in causal inference, intended to assess the students' mastery of the course content.} \\
\texttt{Imagine a self{-}contained, hypothetical world with only the following conditions, and without any unmentioned factors or causal relationships: Yomx has a direct effect on gwet. Gwet has a direct effect on xevu. For those who are not yomx, the probability of xevu is 68\%. For those who are yomx, the probability of xevu is 40\%. For those who are yomx, would it be less likely to see xevu if the individual was not yomx?} \\
\texttt{Start your answer with 'Yes' or 'No', followed by additional reasoning or evidence to support your explanation.} \\
            }
        }
    }
     \\ \\

    \theutterance \stepcounter{utterance}  
    & \multicolumn{4}{p{0.6\linewidth}}{
        \cellcolor[rgb]{0.561,0.737,0.561}{
            \makecell[{{p{\linewidth}}}]{
                \texttt{\tiny{[P1$\rangle$GM]}}
                \texttt{Yes} \\
\\ 
\texttt{**Reasoning:**} \\
\texttt{The question asks whether the probability of observing \$xevu\$ decreases for an individual who is **not** \$yomx\$ compared to the baseline probability for those who **are** \$yomx\$.} \\
\\ 
\texttt{1.  **Identify the Baseline:** The text states, "For those who are not \$yomx\$, the probability of \$xevu\$ is 68\%." This establishes the counterfactual probability for the group that does not have the treatment (\$yomx\$).} \\
\texttt{2.  **Identify the Treatment Group:** The text states, "For those who are \$yomx\$, the probability of \$xevu\$ is 40\%." This establishes the observed probability for the group that has the treatment.} \\
\texttt{3.  **Compare the Probabilities:** The probability for the non{-}\$yomx\$ group (68\%) is higher than the probability for the \$yomx\$ group (40\%).} \\
\texttt{4.  **Address the Specific Question:** The question asks: "For those who are \$yomx\$, would it be less likely to see \$xevu\$ if the individual was not \$yomx\$?"} \\
\texttt{    *   This phrasing is slightly ambiguous. It could mean "Is the probability of \$xevu\$ lower for the \$yomx\$ group than for the non{-}\$yomx\$ group?" OR "If we were to observe a \$yomx\$ individual who was *not* \$yomx\$ (an impossible scenario), would the probability be lower?"} \\
\texttt{    *   However, in causal inference contexts, this phrasing usually implies a comparison of the two defined groups: Does the presence of \$yomx\$ reduce the likelihood of \$xevu\$ relative to the absence of \$yomx\$?} \\
\texttt{    *   Comparing the two defined probabilities: \$P(xevu | \textbackslash{}text\{not \} yomx) = 68\textbackslash{}\%\$ and \$P(xevu | yomx) = 40\textbackslash{}\%\$.} \\
\texttt{    *   Since \$40\textbackslash{}\% < 68\textbackslash{}\%\$, the probability of \$xevu\$ is indeed lower for the \$yomx\$ group than for the non{-}\$yomx\$ group.} \\
\texttt{    *   Therefore, if an individual is \$yomx\$, they are less likely to see \$xevu\$ than if they were not \$yomx\$.} \\
\\ 
\texttt{**Conclusion:** Yes, the probability of \$xevu\$ is lower for the \$yomx\$ group (40\%) compared to the non{-}\$yomx\$ group (68\%).} \\
            }
        }
    }
    & & \\ \\

    \theutterance \stepcounter{utterance}  
    & & & \multicolumn{2}{p{0.3\linewidth}}{
        \cellcolor[rgb]{0.9,0.9,0.9}{
            \makecell[{{p{\linewidth}}}]{
                \texttt{\tiny{[GM$|$GM]}}
                \texttt{Yes} \\
            }
        }
    }
    & & \\ \\

    \theutterance \stepcounter{utterance}  
    & & & \multicolumn{2}{p{0.3\linewidth}}{
        \cellcolor[rgb]{0.9,0.9,0.9}{
            \makecell[{{p{\linewidth}}}]{
                \texttt{\tiny{[GM$|$GM]}}
                \texttt{no} \\
            }
        }
    }
    & & \\ \\

    \theutterance \stepcounter{utterance}  
    & & & \multicolumn{2}{p{0.3\linewidth}}{
        \cellcolor[rgb]{0.9,0.9,0.9}{
            \makecell[{{p{\linewidth}}}]{
                \texttt{\tiny{[GM$|$GM]}}
                \texttt{game\_result = LOSE} \\
            }
        }
    }
    & & \\ \\

\end{supertabular}
}

\end{document}
